\section{Experimental Setup}
\label{sec:setup}

\subsection{Datasets and Cross-Dataset Protocol}
\label{sec:setup:datasets}

We evaluate all supervised methods under a source-to-target
cross-dataset protocol.
The source collection is the HarMeme family, which combines the
COVID-19 and U.S.\ politics domains while retaining their original
domain identities.
The COVID-19 domain contains 3,544 memes and the politics domain
contains 3,469 memes, resulting in 7,013 source samples.
The normalized Facebook Hateful Memes collection, hereafter
\emph{FHM}, contains 9,000 samples and is used exclusively as a
held-out target benchmark.
Table~\ref{tab:dataset_protocol} summarizes the dataset roles.

\begin{table}[t]
    \centering
    \small
    \caption{Dataset roles in the cross-dataset protocol.
    FHM is excluded from every development decision and is used
    only for final held-out evaluation.}
    \label{tab:dataset_protocol}
    \begin{tabular}{lllr}
        \toprule
        Dataset & Family & Role & Samples \\
        \midrule
        Harm-C   & HarMeme & Train/validation & 3,544 \\
        Harm-P   & HarMeme & Train/validation & 3,469 \\
        FHM      & FHM     & Held-out test    & 9,000 \\
        \bottomrule
    \end{tabular}
\end{table}

We construct a single immutable source split using approximately
80\% of HarMeme for training and 20\% for validation.
The split is stratified jointly by the original source domain
(Harm-C or Harm-P) and the binary harmfulness label.
The split seed is fixed to 42 and is independent of the model
initialization seeds.
Every model, baseline, ablation, and knowledge condition uses the
same source split.
The resulting sample assignments and their cryptographic hashes are
stored in persistent split manifests.

FHM is not used for training, validation, early stopping,
checkpoint selection, hyperparameter search, threshold selection,
calibration, prompt development, few-shot demonstration selection,
retrieval-index construction, or knowledge-source selection.
For retrieval-based methods, the labeled meme index is constructed
from the HarMeme training partition only.
All model and decoding decisions are frozen before inference on
FHM.
Memotion is retained in the repository for future experiments but
is disabled in the current protocol because its unified
harmfulness labels require additional validation.

% TODO_CITATION: Add the authoritative HarMeme and FHM citations
% using the final BibTeX keys in reference.bib.

\subsection{Annotation and Supervision}
\label{sec:setup:annotations}

For harmfulness detection, HarMeme and FHM use their original
binary dataset labels as normalized harmful and non-harmful
supervision.
Structured target, intent, and tactic fields are obtained from the
existing schema-validated normalized annotations.
For source training, we use the \texttt{clean} normalized label
set and retain only labels that satisfy the field-specific validity
and confidence policy.

The supervised structured fields are target presence, target
granularity, primary intent, rhetorical tactic, and multimodal
relation.
Harmfulness, target presence, target granularity, primary intent,
and multimodal relation are single-label tasks.
Rhetorical tactic prediction is multi-label.
For every field, labels marked as ignored, unknown, or otherwise
invalid under the canonical vocabulary are excluded through an
explicit supervision mask rather than being treated as negative
instances.

The FHM structured annotations are agent-generated,
schema-validated silver annotations.
Accordingly, we report FHM target, intent, and tactic results as
\emph{silver structured evaluation} and do not describe these
labels as human-adjudicated gold.
For transparency, each structured metric is accompanied by the
number and proportion of FHM samples with valid supervision.
The original annotation payloads, normalized labels, audit flags,
confidence metadata, and label provenance remain available in the
saved prediction artifacts.

\subsection{Compared Methods}
\label{sec:setup:baselines}

We organize compared methods into three groups.

\paragraph{Unimodal and simple multimodal baselines.}
The text-only baseline encodes the OCR text with
DeBERTa-v3-base and trains a task classifier over the resulting
representation.
The image-only baseline uses the local pretrained OpenCLIP
ViT-B/32 image encoder.
The image--text concatenation baseline independently encodes the
image and OCR text, concatenates the two representations, and
applies a multilayer classifier.
We additionally consider a supervised OpenCLIP classifier that
jointly uses the image and text representations without the
evidence extraction, knowledge verification, or task-aware
reasoning components of our framework.

\paragraph{Supervised multimodal methods.}
Our experiment registry prioritizes representative harmful meme
detectors based on CLIP interaction, retrieval-guided contrastive
learning, target-aware modeling, and prompt-based implicit
knowledge.
These include HateCLIPper, RGCL, MOMENTA, and PromptHate.
Reasoning- and explanation-oriented methods such as IntMeme,
SAFE-MEME, and EXPLAINHM++ are included when an official
implementation, compatible checkpoint, and reproducible adapter are
available.
For external methods, we preserve the original architecture and
training objective and modify only the dataset input, output
serialization, and metric interfaces required by the common
HarMeme-to-FHM protocol.
Only completed runs that pass the common artifact and leakage audit
are included in final reproduced-result tables.
Literature-reported scores are stored separately from scores
reproduced under our protocol.

\paragraph{LMM, VLM, and agentic baselines.}
We register direct, chain-of-thought, retrieval-augmented, and
debate-style variants of general-purpose multimodal models.
Their prompts, demonstrations, output schemas, and decoding
policies are developed using HarMeme only.
No FHM example is used as an in-context demonstration or for
prompt selection.
API- or checkpoint-dependent methods are reported only when the
exact model version, inference configuration, and execution cost
are fixed and recorded.

\paragraph{Our model and ablations.}
We compare the complete framework against four train-time
ablations:
without external retrieval,
without task-specific support verification,
without task-aware gating, and
without structured auxiliary supervision.
The same ablation condition is applied during training, validation,
checkpoint selection, and FHM inference.
The structured-auxiliary ablation removes the target-presence and
multimodal-relation losses while retaining compatible output
serialization.

% TODO_CITATION: Insert citations for every external baseline that
% remains in the final experiment set.
% TODO_RESULT: Remove models that are not successfully reproduced.

\subsection{Implementation and Model Selection}
\label{sec:setup:implementation}

Our implementation uses a locally stored pretrained OpenCLIP
ViT-B/32 visual encoder and a locally stored
DeBERTa-v3-base text encoder.
All pretrained assets are verified before training, and their
paths, loading states, and SHA-256 hashes are recorded in each run
manifest.
OCR text is loaded from the existing dataset resources and is not
re-extracted during model training.

For our framework, all valid structured tasks are optimized jointly
using the masked multi-task objective described in
Section~\ref{sec:problem:objective}.
Single-label outputs use cross-entropy losses, whereas rhetorical
tactic prediction uses a multi-label binary classification loss.
The exact task weights, optimizer settings, learning rate, batch
size, number of epochs, early-stopping patience, and hidden
dimensions are fixed in the experiment configuration and reported
in the supplementary material.

The best checkpoint is selected using harmfulness Macro-F1 on the
HarMeme validation partition.
No FHM prediction is generated during checkpoint selection.
For rhetorical tactic decoding, a global sigmoid threshold is
selected on HarMeme validation predictions using non-\texttt{none}
Macro-F1.
Ties are resolved deterministically according to the predefined
decoding contract.
The selected threshold is then frozen and applied unchanged to
FHM.
No target-domain label is inspected while selecting the
checkpoint or decoding threshold.

We first conduct a one-seed pilot with model seed 42.
The final multi-seed evaluation uses model seeds
$\{42,52,123,777,2026\}$ while retaining the same source split.
Experiments are executed on NVIDIA RTX 3090 GPUs.
Each run records the resolved configuration, software environment,
pretrained asset provenance, runtime, parameter counts, peak GPU
memory, selected checkpoint, and decoding policy.

% TODO_IMPLEMENTATION: Verify the exact final optimizer,
% batch size, epoch count, and patience against the frozen
% experiment configuration before submission.

\subsection{Evaluation Metrics}
\label{sec:setup:metrics}

For harmfulness detection, the primary metric is Macro-F1.
We additionally report accuracy, precision, recall, and AUROC when
a valid binary score and both classes are available.

For target presence, target granularity, primary intent, and
multimodal relation, we report Macro-F1 over valid, non-ignored
samples.
For rhetorical tactic prediction, we report logits-only
multi-label Macro-F1 and Micro-F1.
Per-label F1, exact-match ratio, and the separate behavior of the
\texttt{none} fallback are reported in the supplementary material.
Every structured result includes its valid sample count, coverage,
unknown count, ambiguous count, and masked count.

For core ablations, we report both predictive quality and
complexity.
The latter includes total and trainable parameters, training time,
held-out inference latency, peak GPU memory, average retrieved
candidate count, and average verified knowledge count.

For multi-seed experiments, results are summarized as mean and
standard deviation.
Paired significance tests are performed only for runs sharing the
same immutable source and target manifests and the same seed set.
We do not report statistical significance from one-seed pilot
experiments.

% TODO_STATISTICS: Insert the final paired-test method and
% multiple-comparison policy after the five-seed protocol is frozen.

\subsection{Evidence, Rationale, and Verifier Assessment}
\label{sec:setup:explanation_evaluation}

We distinguish automatically measurable diagnostic properties from
human judgments of explanation quality.
Automatic diagnostics include the numbers of selected internal and
external evidence items, candidate-origin distributions, accepted
and rejected knowledge counts, evidence--rationale lexical overlap,
rationale length, provenance completeness, and rule-based
unsupported-claim indicators.
These measures are treated as analysis proxies and not as direct
proof of faithfulness.

The evaluation pipeline additionally exports human-review packages
for evidence relevance and sufficiency, rationale support and
grounding, unsupported claims, and Stage C relevance, validity, and
task-specific support.
Human-evaluation results are included only after valid annotations
have been imported and agreement checks have been completed.
Until then, the corresponding registry entries remain explicitly
marked as blocked by missing human annotation.

\subsection{Reproducibility and Artifact Audit}
\label{sec:setup:reproducibility}

Every paper-targeted run produces a resolved configuration,
environment snapshot, training log, runtime and complexity record,
validation and FHM predictions, metric file, decoding-threshold
record, and audit report.
Run manifests include hashes of the source split, target manifest,
label vocabulary, normalized annotation snapshot, configuration,
source code, and pretrained assets.

A strict audit verifies that the expected losses, prediction fields,
structured logits, decoding artifacts, model checkpoints, and
provenance records are present.
It also verifies that FHM does not occur in source training,
validation, threshold selection, few-shot demonstrations, or
retrieval indexes.
A run is considered complete for paper reporting only when all
required artifacts are present and the strict audit passes.